% ---------------------------------------------------------------------
% Section 7
% ---------------------------------------------------------------------

\begin{cnread}
上一节的 $\Lambda(\square)$ 是个「闭回路可观测量」，但它只是和乐，还不是一个
\emph{联络}。\AEG{} 手稿里的原始提法是：要一个规范等变的比较量，限制到
Borel 扇区 $B_\infty$ 时\textbf{退化回} Paper~0 的仿射平移挠率，而在闭回路上
又\textbf{不能被望远镜抵消}。本节把这两个要求在同一个对象里同时满足：把每条
历史边的算子步提升为半直积群 $\mathrm{Aff}(1,K)=K^\times\ltimes K$ 中的
\textbf{射影挠率传递元} $\mathcal T(e)=(\Lambda(e),\Delta\hat\xi(e))$。两个分量
一个管乘性曲率、一个管加性挠率，限制到 $b\equiv0$ 时前者归 $1$、后者一字不差
地退回 $\hat\xi_n$。本节同时指出附带的规范化因子 $2$ 问题（附录 B）。
\end{cnread}

\subsection{The problem, as \AEG{} states it}

Paper~IV, Section~13, lists as its first open problem:

\begin{quote}\small
\textbf{Projective torsion.}  Construct a gauge-equivariant comparison on
$G\to G/H$ that reduces to affine translation torsion in the Borel sector and
has a nontrivial closed-loop observable not forced to telescope \cite{AEG4}.
\end{quote}

Paper~IV, Proposition~4.9, explains why the naive attempt fails: if the edge
label of a lifted step is defined as the ratio $r_t=h_th_{t-1}^{-1}$ of vertex
frame lifts, then along every closed lifted path
$r_tr_{t-1}\cdots r_{s+1}=h_th_s^{-1}$ telescopes pairwise, so every closed path
has trivial product.  A nontrivial holonomy therefore requires the edge datum to
couple the base motion to the previous frame, not just to record the frames.

\subsection{The incidence-coupled transport}

The needed coupling is supplied by the projectors of
Theorem~\ref{thm:projector} without any additional structure: the natural pairing
between a new image point and the preceding kernel covector is
$\varphi_{b_{k-1}}(v_{a_k})=1+a_kb_{k-1}$.

\begin{definition}[Projective torsion transport,
{\cite[Def.~4.10]{AEG4}}]
\label{def:torsion-transport}
Let $K$ be a field, let $e_k:g_{k-1}\to g_k$ be a lifted step between regular
bivaluations with frame coordinates
\begin{equation}\label{eq:frame-lift}
  g_k=\begin{pmatrix}1&-a_k\\-b_k&1\end{pmatrix}
      \begin{pmatrix}\Phi_k&0\\0&1\end{pmatrix},
  \qquad \Phi_k\in K^{\times}\cong H,
\end{equation}
and suppose the step is transverse, $1+a_kb_{k-1}\ne0$.  The projective torsion
transport of $e_k$ is the element of the affine semidirect product
$\Aff(1,K)=K^{\times}\ltimes K$ given by
\begin{equation}\label{eq:torsion-transport}
  \mathcal T(e_k):=\begin{pmatrix}\Lambda_k&\Delta\hat\xi_k\\[2pt]0&1\end{pmatrix}
  \in\Aff(1,K),
\end{equation}
with multiplicative and additive components
\begin{equation}\label{eq:torsion-components}
  \Lambda_k=\frac{\Phi_k}{\Phi_{k-1}}\cdot\frac{1+a_kb_{k-1}}{1+a_{k-1}b_{k-1}}
  \in K^{\times},
  \qquad
  \Delta\hat\xi_k=-\Phi_{k-1}^{-1}\,\frac{a_k-a_{k-1}}{1+a_{k-1}b_{k-1}}\in K .
\end{equation}
The transport of a history $\gamma=(s_1,\dots,s_n)$ is the chronological
product $\mathcal T(\gamma)=\mathcal T(e_n)\cdots\mathcal T(e_1)$.
\end{definition}

Two remarks on the definition.  The first component generalises the base factor
\eqref{eq:lambda-edge} of Section~6 by multiplying it with the scale ratio of the
step; the second projects the image displacement onto the previous kernel
covector and normalises by the frame scale.  Both are forced once one asks that
the pair $(\Lambda,\Delta\hat\xi)$ be multiplicative along a history and
$\Aff(1,K)$-valued.

\subsection{The Borel reduction}

\begin{theorem}[Two checks, {\cite[Thm.~4.11]{AEG4}}]
\label{thm:torsion}
Let $\mathcal T$ be as in Definition~\ref{def:torsion-transport}.
\begin{enumerate}[label=(\alph*)]
  \item \emph{(Exact Borel reduction.)}  On the affine sector
  $B_\infty$, where $b_k\equiv0$ and $a_k=-\xi_k$ for every prefix
  $g_k(z)=\Phi_kz+\xi_k$, the base correction is identically $1$ and
  \begin{equation}\label{eq:borel-reduction}
    \Lambda_k\big|_{B_\infty}=\frac{\Phi_k}{\Phi_{k-1}}=s_k,
    \qquad
    \Delta\hat\xi_k\big|_{B_\infty}=\Phi_{k-1}^{-1}(\xi_k-\xi_{k-1}).
  \end{equation}
  Multiplying the matrices chronologically gives upper-left entry $\Phi_n$ and
  upper-right entry $\Phi_n\sum_k\Phi_k^{-1}(\xi_k-\xi_{k-1})=\xi_n$, which is
  the affine cocycle law of Paper~0, Proposition~7.1.  In particular, for two
  scale-compatible histories the translation component of
  $\mathcal T(\delta)^{-1}\mathcal T(\gamma)$ is exactly the relative affine
  defect $\Phi^{-1}(\xi_\gamma-\xi_\delta)$ of Paper~0, Section~7.4.
  \item \emph{(Non-telescoping holonomy off the sector.)}  On the four-step loop
  $\square$ of Theorem~\ref{thm:cross-ratio} with closed frame lift
  $\Phi_4=\Phi_0$, the multiplicative component of $\mathcal T$ reproduces
  $\Lambda(\square)$ of \eqref{eq:cross-ratio}.  Hence the transport has a
  closed-loop observable that is gauge-invariant and is not forced to telescope.
  \item \emph{(Modulated additive torsion.)}  Off $B_\infty$ the additive
  component carries the dynamic pole in its denominator, so the accumulated
  translation is a projective torsion modulated by the moving ripple pole rather
  than the affine one.
\end{enumerate}
\end{theorem}

\begin{proof}
(a) Put $b_k=0$ and $a_k=-\xi_k$ in \eqref{eq:torsion-components}: the factor
$(1+a_k\cdot0)/(1+a_{k-1}\cdot0)$ is $1$, leaving
$\Lambda_k=\Phi_k/\Phi_{k-1}$ and
$\Delta\hat\xi_k=-\Phi_{k-1}^{-1}(-\xi_k+\xi_{k-1})
=\Phi_{k-1}^{-1}(\xi_k-\xi_{k-1})$.  The matrix product of the
$\Aff(1,K)$-elements is
$\begin{pmatrix}\Phi_n&\Phi_n\sum\Phi_k^{-1}(\xi_k-\xi_{k-1})\\0&1\end{pmatrix}$,
and the sum telescopes to $\xi_n$ because $\xi_0=0$ (Paper~0,
Proposition~7.1).  (b) is Theorem~\ref{thm:cross-ratio}, whose gauge invariance
was proved there.  (c) is the presence of $1+a_{k-1}b_{k-1}$ in the denominator
of $\Delta\hat\xi_k$, which is not a function of the charges alone once
$b\ne0$.
\end{proof}

\begin{remark}[Why the two demands can be met at once]
\label{rem:two-demands}
The Borel reduction of part (a) and the non-telescoping holonomy of part (b) look
like competing requirements, because a quantity that is multiplicative and
vanishes on $b\equiv0$ might be expected to vanish identically.  They coexist
because they live in different components of $\Aff(1,K)$: the reduction uses the
additive component (the translation $\Delta\hat\xi$, which becomes the affine
cocycle), while the obstruction uses the multiplicative component (the cross
ratio).  On the Borel sector the multiplicative component degenerates to the
step's own scale ratio, which telescopes as it must; off it, the same component
carries the extra base factor and does not.
\end{remark}

\subsection{The curvature of the transport}

For $K=\R$ or $\CC$ one can pass to the infinitesimal picture.  Paper~IV,
Remark~4.12, records the reductive connection whose solder component and
connection one-form are
\begin{equation}\label{eq:connection}
  \omega_H=\Phi^{-1}\mathrm d\Phi+\frac{a\,\mathrm db-b\,\mathrm da}{1+ab},
  \qquad
  \theta_+=\frac{\Phi^{-1}\mathrm da}{1+ab},
\end{equation}
and it records two facts about them that were stated loosely in an earlier
version of that remark and are now exact.  First, its curvature is
$2(1+ab)^{-2}\mathrm da\wedge\mathrm db$, twice the para-K\"ahler form; the
first-order transport density of the cross-ratio holonomy is that form with a
minus sign, so the transport is the holonomy of the connection whose transverse
part is halved (Paper~IV, Proposition~4.17).  Second, on the Borel sector the
solder component evaluates to
\begin{equation}\label{eq:solder-borel}
  \theta_+\big|_{b=0}=\Phi^{-1}\mathrm da=-e^{-M}\mathrm d\xi ,
\end{equation}
because the image coordinate satisfies $a=-\xi$ there; this is \emph{minus} the
left Maurer--Cartan translation form $e^{-M}\mathrm d\xi$ of Paper~0,
Proposition~7.4, Equation~(60).  The minus sign is the sign of $a=-\xi$, not a
different convention for the connection.

The comparison of the two densities is Lemma~\ref{lem:factor-two}.

\begin{lemma}[Two densities, one factor]
\label{lem:factor-two}
The curvature of the connection \eqref{eq:connection} is
$\mathrm d\omega_H=2(1+ab)^{-2}\mathrm da\wedge\mathrm db$, whereas the
first-order log-holonomy density of Corollary~\ref{cor:larmor} is
$-(1+ab)^{-2}\mathrm da\wedge\mathrm db$.  The two differ by the factor $-2$ and
by the choice of orientation.
\end{lemma}

\begin{proof}
Since the $\Phi$-term of $\omega_H$ is closed, only
$\omega_2=(1+ab)^{-1}(a\,\mathrm db-b\,\mathrm da)$ contributes.  With
$f=(1+ab)^{-1}$ one has $\mathrm d\omega_2=\bigl(\partial_a(fa)+\partial_b(fb)
\bigr)\mathrm da\wedge\mathrm db$, and
$\partial_a(fa)+\partial_b(fb)=-2ab(1+ab)^{-2}+2(1+ab)^{-1}=2(1+ab)^{-2}$.  The
transport density is Corollary~\ref{cor:larmor}.  Check~V4.1 verifies the
coefficient symbolically.
\end{proof}

\begin{remark}[Status of the factor]
\label{rem:factor-two}
Lemma~\ref{lem:factor-two} is a normalisation observation, not a contradiction:
the cross-ratio formula \eqref{eq:cross-ratio} is exact and verified, and the
curvature statement of the manuscript is exact as a statement about
\eqref{eq:connection}.  The two normalisations differ by $2$ because the
transport of Definition~\ref{def:torsion-transport} is the holonomy of
$\tfrac12$ of the connection's solder--connection split, i.e. of a connection
whose second term carries a factor $\tfrac12$.  Any statement that identifies
the radiated power with a curvature integral must therefore fix this
normalisation first.  Appendix~\ref{app:curvature} records the computation.

A second observation, unrelated to the factor: three cross-references inside
Paper~IV, Subsection~4.10--4.13 (\texttt{ex:f7-counts},
\texttt{prop:telescoping}, \texttt{eq:chart-projector}) resolved to
\texttt{??} in the build of 20 October 2026, because those labels did not exist.
They have since been repaired in the manuscript source --- the example of
Subsection~3.6 is now labelled, and the other two references point to
Proposition~4.9 and to the projector display \eqref{eq:projector} --- and the
rebuilt Paper~IV reports no undefined reference.
Appendix~\ref{app:register} records the repair together with a second, smaller
one in Paper~0.
\end{remark}

\subsection{The finite witness, restated}

Example~\ref{ex:f7} applies verbatim to the transport: on
$K=\F_7$ with the loop $\square$ taking
$(0,0)\to(1,0)\to(1,1)\to(0,1)\to(0,0)$ and $\Phi_k\equiv1$, one gets
$\Lambda(\square)=4=2^{-1}\in\F_7^{\times}$, a nonidentity element of
$H$.  With $\Phi_k\equiv1$ the additive components are
$\Delta\hat\xi_k=-(a_k-a_{k-1})/(1+a_{k-1}b_{k-1})$, so the transport is a purely
finite, exactly computable object.  Appendix~\ref{app:matrices} lists the four
edge transports of this loop explicitly.

\subsection{Where this leaves the open problem}

Theorem~\ref{thm:torsion} supplies both demands of the problem statement: the
transport reduces exactly to affine torsion on $B_\infty$, and off $B_\infty$ it
has a closed-loop observable that is gauge-invariant and provably not forced to
telescope, with a finite witness.  What the theorem does not supply is a
physical calibration of the transport's \emph{units}: no statement is made about
which physical quantity $\Lambda(\square)-1$ measures, beyond the dictionary of
Hypothesis~\ref{hyp:radiation}.  That is the open end recorded in
Section~\ref{sec:summary}.
